\documentclass[a4paper,11pt]{article}
\usepackage[utf8]{inputenc}
\usepackage[T1]{fontenc}
\usepackage[french]{babel}
\usepackage[left=2cm,right=2cm,top=2cm,bottom=2cm]{geometry}
\usepackage{xcolor}
\usepackage{hyperref}
\usepackage{fontawesome5}
\usepackage{parskip}

% --- Couleurs Air France ---
\definecolor{primary}{RGB}{0, 63, 135}
\hypersetup{colorlinks=true, urlcolor=primary}

\begin{document}

% --- EXPÉDITEUR ---
\noindent
\begin{minipage}[t]{0.5\textwidth}
    \textbf{Loïc Aron MBASSI EWOLO}\\
    Élève-Ingénieur Bac+4 -- Double diplôme ENSEA\\[3pt]
    \faMapMarker\ Cergy, France\\
    \faPhone\ (+33) 7 59 52 33 98\\
    \faEnvelope\ \href{mailto:loicaron.mbassiewolo@ensea.fr}{loicaron.mbassiewolo@ensea.fr}
\end{minipage}
\hfill
% --- DATE ---
\begin{minipage}[t]{0.4\textwidth}
    \raggedleft
    Cergy, le 6 février 2026
\end{minipage}

\vspace{1.2cm}

% --- DESTINATAIRE ---
\hfill
\begin{minipage}[t]{0.5\textwidth}
    \raggedleft
    \textbf{Air France}\\
    Direction de l'Ingénierie Cabine\\
    45 rue de Paris\\
    95747 Roissy CDG Cedex
\end{minipage}

\vspace{1cm}

\noindent\textbf{Objet :} Candidature au stage Développeur banc de test graphique 2D (Réf. 2026-23813)

\vspace{0.8cm}

Madame, Monsieur,

Le parsing XML et la conception d'interfaces graphiques interactives constituent le cœur de mon projet open source Laravel API Generator, où j'analyse géométriquement des fichiers draw.io pour générer du code automatiquement. Cette expérience technique, combinée à ma formation d'ingénieur à l'ENSEA spécialisée en systèmes embarqués, me permet de candidater avec enthousiasme au poste de développeur banc de test graphique 2D au sein du DOA d'Air France.

L'optimisation d'une application Python pour valider la programmation des fonctions cabine de l'A350 rejoint directement mes compétences. Sur mon projet de générateur de code, j'ai conçu un parseur XML avancé capable d'interpréter des structures complexes et de calculer les relations géométriques entre entités. Ce travail m'a permis de maîtriser le traitement de données structurées et la génération d'interfaces utilisables par des non-développeurs. Le projet, aujourd'hui open source avec plus de 30 étoiles sur GitHub, démontre ma capacité à produire des outils robustes et documentés.

Mon expérience en développement d'interfaces graphiques s'étend également à des contextes variés. Pour un projet de serious game médical, j'ai conçu une IHM interactive sous Unity/C\# intégrant des scénarios cliniques avec retours en temps réel. Par ailleurs, lors de mon stage de recherche au MILA cet été, j'ai développé des pipelines Python complets pour l'analyse de données expérimentales, renforçant ma rigueur dans le traitement et la visualisation de résultats. Ces expériences m'ont appris à concevoir des interfaces ergonomiques tout en respectant des contraintes techniques strictes.

L'environnement aéronautique m'attire particulièrement par son exigence de fiabilité et de traçabilité. Contribuer à l'amélioration d'un outil de test pour les systèmes cabine de nouvelle génération représente une opportunité d'appliquer mes compétences Python et XML dans un contexte industriel concret. Je suis disponible dès mars 2026 pour une durée de 4 à 6 mois.

Je reste à votre disposition pour un entretien afin de discuter de ma candidature et de ma motivation pour ce stage.

Je vous prie d'agréer, Madame, Monsieur, l'expression de mes salutations distinguées.

\vspace{0.8cm}

\textbf{Loïc Aron MBASSI EWOLO}\\
\textit{Élève-Ingénieur ENSEA / Polytechnique Yaoundé}

\end{document}
